\documentclass[10pt,a4paper]{amsart}
\usepackage[margin=19mm]{geometry}
\usepackage{amsmath,amssymb,mathtools}
\usepackage[scr=rsfs,cal=euler]{mathalfa}
\usepackage{xfrac}
\usepackage[unicode,hidelinks,bookmarks=false]{hyperref}
\newcommand{\ie}{i.e.\ }
\newcommand{\cf}{cf.\ }
\newcommand{\resp}{respectively\ }
\newcommand{\Subsection}{\subsection}
\providecommand{\nobreakdash}{\nobreak\hbox{-}}
\setlength{\emergencystretch}{2em}
\multlinegap0pt
\usepackage{tikz}
\usepackage{amssymb}
\usepackage{xcolor}
\usepackage{enumerate}
\usepackage[shortlabels]{enumitem}
\newtheorem{thm}{Theorem}
\usepackage{cleveref}
\crefname{thm}{Theorem}{Theorems}
\DeclareMathOperator{\CB}{\mathbf{C}}
\DeclareMathOperator{\Irr}{Irr}
\title{Subnormalizers and character correspondences in $p$-solvable groups}
\author{Gabriel A. L. Souza}
\date{Source: arXiv:2605.22669v1 (CC BY 4.0)}
\begin{document}
\maketitle

\noindent\emph{Source and licence.} Subnormalizers and character correspondences in $p$-solvable groups. Version arXiv:2605.22669v1 (CC BY 4.0), distributed by arXiv under the Creative Commons Attribution 4.0 International licence, http://creativecommons.org/licenses/by/4.0/, as recorded in the arXiv licence metadata for this exact version and shown on the abstract page https://arxiv.org/abs/2605.22669v1. Full licence notice: \texttt{licenses/arxiv-2605.22669-CC-BY-4.0.txt}.\par\smallskip

\noindent\textit{Editorial scope.} Counted unit: the article's thm correspondence (source0.tex lines 357--368). The article's complete original proof of it is delivered with it (source0.tex lines 370--404, one \texttt{proof} environment of 35 lines). No mathematics is written, repaired or re-derived: the statement and the proof are the article's own lines, in the article's own notation, and no retained line is substituted or re-flowed.  No further source-local context is needed: every cross-reference of every retained line resolves inside the two delivered spans.  Nothing was omitted from any retained span, so the package carries no omission record.  No retained line contains a citation, so the wrapper carries no bibliography and no bibliography entry.  No retained line contains a figure environment, an image-inclusion command or drawing code, so no image is delivered and the package declares no asset dependency.  Editorial formatting only, in addition: an AMS \texttt{amsart} preamble that restates the article's own declarations and macros used by the retained lines, namely the environments \texttt{thm} on separate counters, together with the standard packages those lines need.  The retained environments print in this document as Theorem 1 in order of appearance, whereas the article prints the same items with the numbering of its own full document.  The complete native source of the article as distributed in the arXiv e-print for this exact version is bundled unchanged as \texttt{sources/201149/source0.tex}.
\begin{thm}
	\label{correspondence}
	Let $p$ be a prime and let $G$ be a group of $p'$-order. Suppose $A$ is a $p$-group acting on $G$ by automorphisms and let $C = \CB_G(A)$ and $C \leq H \leq G$ an $A$-invariant subgroup of $G$. Then, there exists a natural bijection $f_A: \Irr_{A}(G) \to \Irr_{A}(H)$. In fact, if $\theta \in \Irr_{A}(G)$, then
	\begin{equation*}
		\theta_H = ef_A(\theta) + p\Delta + \Xi,
	\end{equation*}
	where $e \equiv \pm 1 \pmod{p}$ and $\Delta, \Xi$ are characters of $H$ (or zero) such that no constituent of $\Xi$ is $A$-invariant. Also, if $\tau \in \Irr_{A}(H)$, then 
	\begin{equation*}
		\tau^G = d\mu + p\Lambda + \rho,
	\end{equation*}
	where $f_A(\mu) = \tau$, $d \equiv \pm 1 \pmod{p}$ and $\Lambda, \rho$ are characters of $G$ (or zero) such that no constituent of $\rho$ is $A$-invariant. Moreover, writing $\varphi = f_A(\theta)$, then, if $*$ denotes the $A$-Glauberman correspondence, $\theta^* = \varphi^*$ and $[\theta_C, \theta^*] \equiv [\theta_H, \varphi][\varphi_C, \theta^*] \pmod{p}$.
\end{thm}

\begin{proof}
	First, there is, indeed a bijection between $\Irr_{A}(G)$ and $\Irr_{A}(H)$ by the Glauberman correspondence applied twice, since $\CB_H(A) = \CB_G(A)$, as $\CB_G(A) \subseteq H$. That is to say, if $\theta \in \Irr_{A}(G)$, then there exists a unique $\varphi \in \Irr_{A}(H)$ such that $\theta^* = \varphi^*$. We will show that $\varphi$ is under the conditions outlined above.
	
	Write $\theta_H = \sum_{\eta \in \Irr(H)} a_\eta \eta$. Since $H$ is $A$-invariant, $A$ acts on $\Irr(H)$, so that we can write $\Irr(H) = \mathcal{O}_1 \sqcup \cdots \sqcup \mathcal{O}_k$ partitioned by $A$-orbits. Also, $\theta$ is $A$-invariant, meaning $[\theta_H, \eta] = [\theta_H, \eta^x]$ for all $x \in A$. Thus, the $a_\eta$ are, in fact, constant on each orbit. Let $\eta_i$ be a representative of $\mathcal{O}_i$ for each $i$ and let $a_i = [\theta_H, \eta_i]$. Then, $\theta_H = \sum_{i = 1}^k a_i \sum_{\eta \in \mathcal{O}_i} \eta$. Notice that, given $\eta \in \Irr(H)$, $(\eta^x)_{C} = \eta_{C}$ for all $x \in A$. Consequently, for all $\eta \in \mathcal{O}_i$, $\eta_{C} = {\eta_i}_{C}$.
	
	Write $\theta_{C} = m\theta^* + p\Delta'$ as in the $A$-Glauberman correspondence. Then, we have
	\begin{equation}
		\label{eq1}
		m\theta^* + p\Delta' = \theta_{C} = \sum_{i = 1}^k a_i \sum_{\eta \in \mathcal{O}_i} \eta_{C} = \sum_{i = 1}^k |\mathcal{O}_i| a_i {\eta_i}_{C}.
	\end{equation}
	From the fact that $m \not \equiv 0 \pmod{p}$, it follows that there is at least one of the $\mathcal{O}_i$ of size $1$ (since $A$ is a $p$-group, they are all powers of $p$). Without loss of generality, we may assume that orbits $1$ through $l$ are the ones of size $1$ (where $1 \leq l \leq k$). Then, \Cref{eq1} becomes
	\begin{equation*}
		m\theta^* + p\Delta' = \sum_{i = 1}^l a_i {\eta_i}_{C} + p \tilde{\Delta},
	\end{equation*}
	for some (possibly zero) character $\tilde{\Delta}$. For $1 \leq i \leq l$, we may write ${\eta_i}_{C} = m_i\eta_i^* + p\Delta_i$ as in the $A$-Glauberman correspondence. Then, we obtain
	\begin{equation}
		\label{eq2}
		m\theta^* + p\Delta' = \sum_{i = 1}^l a_i m_i\eta_i^*+ p \Delta'',
	\end{equation}
	where, again, $\Delta''$ is a character or $0$. Since the $\eta_i$ are pairwise distinct, so are the $\eta_i^*$. Thus, by \Cref{eq2}, there exists exactly one $i_0$ between $1$ and $l$ such that $\eta_{i_0}^* = \theta^*$ and we must also have $a_{i_0} m_{i_0} \equiv m \equiv \pm 1 \pmod{p}$. In particular, since $m_{i_0} \equiv \pm 1 \pmod{p}$, we have that $a_{i_0} \equiv \pm 1 \pmod{p}$. If $i \neq i_0$ is between $1$ and $l$, then $a_i \equiv 0 \pmod{p}$, again by \Cref{eq2}. And, if $l < i \leq k$, then $\eta_i$ is not $A$-invariant. This gives us the first expression in the statement of the theorem.
	
	For the other direction, let $\tau \in \Irr_{A}(H)$ and let $\mu \in \Irr_{A}(G)$ be such that $f_A(\mu) = \tau$. Then, we know that
	\begin{equation*}
		\mu_H = d\tau + p\Delta + \Xi
	\end{equation*}
	for some $d \equiv \pm 1 \pmod{p}$, and some characters $\Delta, \Xi$, the latter of which does not have $\tau$ as a constituent. Then,
	\begin{equation*}
		\tau^G = d\mu + p\Lambda + \rho
	\end{equation*}
	for some characters $\Lambda, \rho$ (or zero), by Frobenius reciprocity. We may assume that, for all constituents $\xi$ of $\rho$, $[\tau^G, \xi] \not \equiv 0 \pmod{p}$. Suppose there exists one such constituent $\xi$ which is $A$-invariant. Then,
	\begin{equation*}
		\xi_H = mf_A(\xi) + p\Phi + \Psi,
	\end{equation*}
	using the direction we already proved. Since $\xi \neq \mu$ and $f_A$ is bijective, $f_A(\xi) \neq \tau$. Also, $\tau$ is $A$-invariant. Then, by Frobenius reciprocity, since $0 \neq [\tau^G, \xi] = [\tau, \xi_H]$, we must have $[\tau, \xi_H] \equiv 0 \pmod{p}$, a contradiction.
\end{proof}

\end{document}
